% 03_diagnostic.tex  --  Background + The Split-Landscape Discrepancy

\section{Background and the SLD Diagnostic}
\label{sec:diagnostic}

\subsection{GBDT Split-Gain Landscape}

Let the full training dataset be $D = \{(x_i, y_i)\}_{i=1}^{N}$ with features
$x_i \in \mathbb{R}^F$ and labels $y_i$. A gradient-boosted decision tree
(GBDT) fits an additive ensemble of trees by performing, at each node, a
greedy split that maximises the regularised gain~\cite{friedman2001gbm,
chen2016xgboost}. Given first- and second-order loss statistics
$g_i = \partial_{\hat{y}} \ell(\hat{y}_i, y_i)$ and
$h_i = \partial^2_{\hat{y}} \ell(\hat{y}_i, y_i)$ for each instance $i$, the
gain of partitioning the instance set $I$ at a node into left and right
children $I_L, I_R$ is
%
\begin{equation}
  \label{eq:gain}
  \begin{aligned}
  \mathcal{G}(I_L, I_R) = \frac{1}{2}\Bigg[
      &\frac{\bigl(\sum_{i \in I_L} g_i\bigr)^2}{\sum_{i \in I_L} h_i + \lambda}
    + \frac{\bigl(\sum_{i \in I_R} g_i\bigr)^2}{\sum_{i \in I_R} h_i + \lambda} \\
    &- \frac{\bigl(\sum_{i \in I} g_i\bigr)^2}{\sum_{i \in I} h_i + \lambda}
    \Bigg] - \gamma,
  \end{aligned}
\end{equation}
%
where $\lambda \geq 0$ is the $\ell_2$ leaf-weight regulariser and
$\gamma \geq 0$ is the minimum gain threshold~\cite{chen2016xgboost}. Modern
histogram-based GBDTs~\cite{ke2017lightgbm,prokhorenkova2018catboost} discretise
each feature $f$ into $B$ bins; per-bin sums of $(g, h)$ suffice to evaluate
every candidate split.

\begin{lemma}[Histogram sufficiency; standard, e.g.~\cite{ke2017lightgbm,ibragimov2019mvs}]
\label{lem:histogram}
Given the current gradients and Hessians, a histogram GBDT with $B$ bins per
feature has its full per-node split-gain landscape determined entirely by the
$F \times B$ grid of per-bin $(g, h)$ aggregates. No other information from the
training set is required to evaluate \eqref{eq:gain} at every candidate
(feature, threshold) pair.
\end{lemma}

\noindent We define the \emph{split-gain landscape}
$\Gamma_D \in \mathbb{R}^{F \times (B-1)}$ as the gain $\mathcal{G}$ that each
(feature, bin-boundary threshold) candidate would achieve on dataset $D$ under a
fixed prediction state. Lemma~\ref{lem:histogram} is a standard consequence of
the histogram GBDT algorithm; we state it here to ground the diagnostic. Its
importance is histogram sufficiency: once the prediction state fixes
$g_i,h_i$, $\Gamma_D$ can be computed in one histogram pass.

\subsection{Split-Landscape Discrepancy}

Let $D'=\{(x'_j,y'_j,w'_j)\}_{j=1}^m$ be a weighted selected or synthetic
surrogate of $D$, with $m \ll N$. Synthetic points are binned using the same
feature preprocessing and bin boundaries as $D$; their labels may be hard labels
or soft labels converted to the objective's expected gradients. We quantify the
structural gap between $D$ and $D'$ via their landscape distance.

\begin{definition}[Split-Landscape Discrepancy]
\label{def:sld}
Let the raw landscape error be
$\delta_p(D,D')=\lVert\Gamma_D-\Gamma_{D'}\rVert_p$. The relative
Split-Landscape Discrepancy of order $p$ between dataset $D$ and surrogate $D'$
is
%
\begin{equation}
  \label{eq:sld}
  \mathrm{SLD}_p(D, D') \;=\;
  \frac{\delta_p(D,D')}
       {\max(\lVert \Gamma_D\rVert_p,\epsilon)},
\end{equation}
%
using the same bins, objective, regularisation, base score, and feature
preprocessing for $D$ and $D'$. We use $p = \infty$ (worst-case relative split
distortion) and the rank-based \emph{split-regret} variant. Root-SLD compares
the root landscape. Reference-tree early-node SLD averages over early nodes of a
fixed full-data reference tree, using the reference node memberships. Full-path
SLD extends the same comparison along deeper reference-tree paths.
\end{definition}

\noindent Both variants are computed from histogram statistics under a fixed
prediction state, and crucially \emph{no model is retrained on $D'$}.
\Cref{fig:pipeline} illustrates how SLD fits into the condensation evaluation
pipeline as a no-candidate-retraining alternative to the retrain-and-evaluate
loop.

\subsection{Diagnostic Sanity Checks}

The diagnostic rests on two standard facts about histogram GBDTs and set
coverage. We use them only to motivate what SLD measures; the contribution of
the paper is the empirical measurement study across 17 condensation methods.

\paragraph{Margin preservation.}
Let $\Delta_v$ be the gain margin between the best-ranked and second-best-ranked
split at node $v$ on $D$. If
%
\begin{equation}
  \label{eq:prop1}
  \delta_\infty(D, D') < \Delta_v / 2,
\end{equation}
%
then the argmax split selected by $D'$ coincides with the argmax selected by $D$
at node $v$, assuming a unique argmax and deterministic tie-breaking. In the
dose-response cells satisfying this raw-error margin condition, root-level split
agreement equals $1.0$ (\cref{sec:results-dose}).

\paragraph{Coverage guarantee.}
The split-coverage objective is the unweighted fraction of full-data high-gain
split thresholds, the top-$k$ full-data split candidates used by the coverage
selector, whose adjacent bins are represented by at least one selected or
weighted surrogate point. It is binary threshold coverage, not gain-weighted
path fidelity. As a standard set-coverage objective, it is monotone submodular
under a cardinality constraint, so greedy maximisation has the usual
$(1 - 1/e)$ guarantee~\cite{nemhauser1978submodular}. This defines
HistDistill-Greedy: it covers the split landscape near-optimally yet performs
poorly by downstream predictive performance
(\cref{sec:results-coverage}). The gap between coverage and predictive fidelity is the
central empirical finding of \cref{sec:results-coverage}.

\subsection{Computing SLD: Algorithm}

\Cref{alg:sld} summarises the SLD computation. The key property is that, after
the fixed prediction state is available, scoring a candidate requires only one
histogram pass and one arithmetic comparison of landscapes; no candidate GBDT is
trained. After shared binning, histogram accumulation costs $O(NF)$ for full
data and $O(mF)$ for the surrogate; prefix-sum gain evaluation and memory are
$O(FB)$ per node.

\begin{algorithm}[!ht]
\caption{Compute Split-Landscape Discrepancy without candidate retraining}
\label{alg:sld}
\begin{algorithmic}[1]
\Require Full dataset $D$, weighted surrogate $D'$, bins, norm order $p$, fixed prediction state $\pi$
\Ensure $\mathrm{SLD}_p(D, D')$ and split-regret score
\State Compute first- and second-order loss statistics for $D$ under $\pi$
       \Comment{root base score or reference-booster predictions}
\State Build histogram $H_D \in \mathbb{R}^{F \times B \times 2}$:
       for each feature $f$ and bin $b$, store
       $\bigl(\textstyle\sum_{i: x_{if} \in b} g_i,\;
               \textstyle\sum_{i: x_{if} \in b} h_i\bigr)$
\State Compute landscape $\Gamma_D \in \mathbb{R}^{F \times (B-1)}$ via \eqref{eq:gain}
       applied to every candidate split
\State Repeat Steps 1--3 for $D'$ to obtain $\Gamma_{D'}$
       \Comment{selected rows reuse weights; synthetic rows are evaluated under $\pi$}
\State $\mathrm{SLD}_p \leftarrow
       \lVert \Gamma_D - \Gamma_{D'} \rVert_p / \max(\lVert\Gamma_D\rVert_p,\epsilon)$
\State \textbf{Split-regret:} let $(f^*, b^*) = \arg\max_{f,b} [\Gamma_D]_{f,b}$
       and $(\hat{f}, \hat{b}) = \arg\max_{f,b} [\Gamma_{D'}]_{f,b}$; \\
       \hspace{4em} split-regret $\leftarrow
         \max\!\left(0,\frac{[\Gamma_D]_{f^*,b^*}-[\Gamma_D]_{\hat f,\hat b}}
         {\max(|[\Gamma_D]_{f^*,b^*}|,\epsilon)}\right)$
\State \Return $\mathrm{SLD}_p$, split-regret
\end{algorithmic}
\end{algorithm}
